Suppose now that \(\omega^{1}_{X_{0}/S_{0}}\) admits a finite presentation
(\cf\ EGA~\(0_{\mathrm{I}}\), 5.2.5), which will in particular be the case if
\(X_{0}\) is \emph{locally of finite presentation} over \(S_{0}\)
(\cf\ EGA~IV\(_4\), 16.4.22). Then, if \(T\) is \emph{flat} over \(S_{0}\), it
follows from EGA~\(0_{\mathrm{I}}\), 6.7.6 that
\[
\mathcal{H}\mathrm{om}_{\mathcal{O}_{T}}
  \bigl(\omega^{1}_{X_{0}/S_{0}}\otimes_{\mathcal{O}_{S_{0}}}\mathcal{O}_{T},\,
  \mathcal{J}\otimes_{\mathcal{O}_{S_{0}}}\mathcal{O}_{T}\bigr)
\simeq
\mathcal{H}\mathrm{om}_{\mathcal{O}_{S_{0}}}
  (\omega^{1}_{X_{0}/S_{0}},\mathcal{J})
  \otimes_{\mathcal{O}_{S_{0}}}\mathcal{O}_{T},
\]
whence
\[
\Hom_{S_{0}}(T,L_{0X})
=
\Gamma\bigl(T,\,
  \mathcal{H}\mathrm{om}_{\mathcal{O}_{S_{0}}}
    (\omega^{1}_{X_{0}/S_{0}},\mathcal{J})
    \otimes_{\mathcal{O}_{S_{0}}}\mathcal{O}_{T}\bigr).
\]
Introducing the notation \(\mathbf{W}(\ )\) of~I, 4.6.1, one has therefore proved
that for every \(S_{0}\)-scheme \(T\) \emph{flat} over \(S_{0}\), one has
\[
\Hom_{S_{0}}(T,L_{0X})
=
\Hom_{S_{0}}\bigl(T,\,
  \mathbf{W}\bigl(\mathcal{H}\mathrm{om}_{\mathcal{O}_{S_{0}}}
    (\omega^{1}_{X_{0}/S_{0}},\mathcal{J})\bigr)\bigr).
\]
In summary, if \(\omega^{1}_{X_{0}/S_{0}}\) \emph{admits a finite presentation},
and if one restricts oneself to the category of \(S\)-schemes \emph{flat}
over~\(S\), one has
\begin{equation*}
(0.6.1)\qquad
L'_{X}
=
\prod_{S_{0}/S}
  \mathbf{W}\bigl(\mathcal{H}\mathrm{om}_{\mathcal{O}_{S_{0}}}
    (\omega^{1}_{X_{0}/S_{0}},\mathcal{J})\bigr),
\end{equation*}
and \(\mathcal{H}\mathrm{om}_{\mathcal{O}_{S_{0}}}
(\omega^{1}_{X_{0}/S_{0}},\mathcal{J})\) is a \emph{quasi-coherent}
\(\mathcal{O}_{S_{0}}\)-module, by EGA~I, 9.1.1.

Let us remark finally that if \(\omega^{1}_{X_{0}/S_{0}}\) is moreover locally
free (of finite rank), for example if \(X_{0}\) is \emph{smooth} over \(S_{0}\)
(in which case it is automatically locally of finite presentation over
\(S_{0}\)), one has
\begin{equation*}
(0.6.2)\qquad
\mathcal{H}\mathrm{om}_{\mathcal{O}_{S_{0}}}
  (\omega^{1}_{X_{0}/S_{0}},\mathcal{J})
\simeq
\uLie(X_{0}/S_{0})\otimes_{\mathcal{O}_{S_{0}}}\mathcal{J},
\end{equation*}
where one denotes by abuse of language (\(X_{0}\) not necessarily being an
\(S_{0}\)-group) \(\uLie(X_{0}/S_{0})\) the dual of the
\(\mathcal{O}_{S_{0}}\)-module \(\omega^{1}_{X_{0}/S_{0}}\).%
{\renewcommand{\thefootnote}{(23)}\nde{This is justified by~II, 3.3 and~4.11.
Indeed, the functor \(L^{\varepsilon_{0}}_{X_{0}/S_{0}}\), ``tangent space to
\(X_{0}\) over \(S_{0}\) at the point \(\varepsilon_{0}\)'', is represented by
the vector fibration
\(\mathbf{V}(\varepsilon_{0}^{*}(\Omega^{1}_{X_{0}/S_{0}}))
=\mathbf{V}(\omega^{1}_{X_{0}/S_{0}})\) over \(S_{0}\), whose sheaf of sections
is the dual module
\(\mathcal{H}\mathrm{om}_{\mathcal{O}_{S_{0}}}
(\omega^{1}_{X_{0}/S_{0}},\mathcal{O}_{S_{0}})\);
the latter is \(\uLie(X_{0}/S_{0})\) when \(X_{0}\) is an \(S_{0}\)-group and
\(\varepsilon_{0}\) the unit section.}}

Proposition~0.2 (and its proof) has two important corollaries.%
{\renewcommand{\thefootnote}{(24)}\nde{One has added to~0.7 the
corollary~0.7.bis, which was used implicitly in the proof of~0.8. Let us
point out here that the numbers~0.8 to~0.12, as well as~0.17, which play an
important technical role in the sequel of this expos\'e, are consequences
of~0.7 and~0.7.bis.}}

\begin{corollary}{0.7}
\label{III.0.7}
Let \(X\) be an \(S\)-scheme.

a)~Every \(S\)-endomorphism of \(X\) inducing the identity on \(X_{\mathcal{J}}\)
is an automorphism.

b)~One has an exact sequence of groups:
\[
0 \to
\Hom_{\mathcal{O}_{X_{0}}}(\Omega^{1}_{X_{0}/S_{0}},\mathcal{J}\mathcal{O}_{X})
\xrightarrow{i}
\Aut_{S}(X)
\to
\Aut_{S_{\mathcal{J}}}(X_{\mathcal{J}}).
\]

c)~Moreover, if one lets \(\Aut_{S}(X)\) operate on the first group by
transport of structure, one has, for all \(u\in\Aut_{S}(X)\) and
\(m\in\Hom_{\mathcal{O}_{X_{0}}}(\Omega^{1}_{X_{0}/S_{0}},
\mathcal{J}\mathcal{O}_{X})\):
\[
i(um)=u\,i(m)\,u^{-1}.
\]
\end{corollary}

\begin{proof}
By~0.2~(i), \(\Hom_{X^{+}}(X,X)\) is a principal homogeneous set under
\(\Hom_{X^{+}}(X,L_{X})\), for it is certainly non-empty; indeed, it contains
a marked point: the identity automorphism of~\(X\).%
{\renewcommand{\thefootnote}{(25)}\nde{One has set out what follows in
detail.}}
Consequently, the map \(m\mapsto m\cdot\mathrm{id}_{X}\) induces a
\emph{bijection}
\[
\Hom_{\mathcal{O}_{X_{0}}}(\Omega^{1}_{X_{0}/S_{0}},\mathcal{J}\mathcal{O}_{X})
=
L_{X}(X)
\isomto
\Hom_{X^{+}}(X,X).
\]
Let \(m\in L_{X}(X)\) and let \(f=m'\cdot\mathrm{id}_{X}\) be an element of
\(\Hom_{X^{+}}(X,X)\). Applying~0.2~(ii) to~\(f\), one obtains that:
\[
f\circ(m\cdot\mathrm{id}_{X})
=
L_{f}(X)(m)\cdot f
=
L_{f}(X)(m)\cdot(m'\cdot\mathrm{id}_{X}).
\]
On the other hand, as \(f\) is an \(X^{+}\)-endomorphism of~\(X\), one has
\(f_{\mathcal{J}}=\mathrm{id}_{X_{\mathcal{J}}}\) and hence
\(f_{0}=\mathrm{id}_{X_{0}}\); as \(L_{f}\) depends only on~\(f_{0}\)
(\cf\ N.D.E.~(17) in~0.1.10), one therefore has \(L_{f}(X)(m)=m\).
Consequently, the equality above is rewritten:
\[
(m'\cdot\mathrm{id}_{X})\circ(m\cdot\mathrm{id}_{X})
=
m\cdot(m'\cdot\mathrm{id}_{X})
=
(m+m')\cdot\mathrm{id}_{X}.
\]
This shows that the bijection \(m\mapsto m\cdot\mathrm{id}_{X}\) transforms
the group law of \(\Hom_{X^{+}}(X,L_{X})\) into the composition law of the
\(X^{+}\)-endomorphisms of~\(X\).

It follows first that every element of \(\Hom_{X^{+}}(X,X)\) is invertible,
which is the first assertion of the statement, and then that one has an exact
sequence
\[
0 \to
\Hom_{X^{+}}(X,L_{X})
\xrightarrow{i}
\Aut_{S}(X)
\to
\Aut_{S_{\mathcal{J}}}(X_{\mathcal{J}}),
\]
which is the second.

Let us remark now that the morphism~\(i\) defined above is functorial in~\(X\)
for isomorphisms, for it is defined in structural terms from the operation
of \(L_{X}\) on~\(X\) over~\(X^{+}\), itself functorial in~\(X\) by
assertion~(ii) of proposition~0.2.%
{\renewcommand{\thefootnote}{(26)}\nde{One has set out what follows in
detail.}}
Therefore every automorphism~\(u\) of~\(X\) over~\(S\) induces by transport of
structure isomorphisms
\[
h\colon\Hom_{X^{+}}(X,L_{X})\isomto\Hom_{X^{+}}(X,L_{X})
\]
and \(f\colon\Aut_{S}(X)\isomto\Aut_{S}(X)\) such that the following diagram is
commutative:
\[
\xymatrix{
\Hom_{X^{+}}(X,L_{X}) \ar[r]^{i} \ar[d]_{h}
  & \Aut_{S}(X) \ar[d]_{f} \\
\Hom_{X^{+}}(X,L_{X}) \ar[r]^{i}
  & \Aut_{S}(X)
}
\]
\ie such that \(f\circ i=i\circ h\). On the other hand, \(f\) is given by the
commutative diagram:
\[
\xymatrix{
X \ar[r]^{a} \ar[d]_{u} & X \ar[d]^{u} \\
X \ar[r]_{f(a)} & X\,,
}
\]
that is, \(f(a)=u\circ a\circ u^{-1}\), for every \(a\in\Aut_{S}(X)\).
Writing \(i(h(m))=f(i(m))\), one finds the desired formula.
\end{proof}

\begin{corollary}{0.7.bis}
\label{III.0.7.bis}
Let \(X\) be an \(S\)-scheme such that \(X_{\mathcal{J}}\) is an
\(S_{\mathcal{J}}\)-monoid. Then \(L_{X}\) is equipped with a structure of
\(S\)-monoid, one has a split exact sequence of \(S\)-monoids:
\[
\xymatrix{
1 \ar[r]
  & L'_{X} \ar[r]^{i}
  & L_{X} \ar@<0.5ex>[r]^{p}
  & X^{+} \ar@<0.5ex>[l]^{s} \ar[r]
  & 1
}
\]
and the monoid law induced on \(L'_{X}\) coincides with its structure of
abelian group.
In particular, if \(X_{\mathcal{J}}\) is an \(S_{\mathcal{J}}\)-group, then
\(L_{X}\) is an \(S\)-group and is the semidirect product of \(X^{+}\) and
of~\(L'_{X}\).
\end{corollary}

\begin{proof}
Indeed, as \(X_{\mathcal{J}}\) is an \(S_{\mathcal{J}}\)-monoid, then
\(X^{+}=\prod_{S_{\mathcal{J}}/S}X_{\mathcal{J}}\) is an \(S\)-monoid (indeed,
one has \(X^{+}(Y)=X_{\mathcal{J}}(Y_{\mathcal{J}})\) for every \(Y\to S\)).
For every \(S\)-scheme~\(Y\), let us denote by \(\widetilde{Y}_{\mathcal{J}}\)
the affine \(Y_{\mathcal{J}}\)-scheme corresponding to the quasi-coherent
\(\mathcal{O}_{Y_{\mathcal{J}}}\)-algebra
\(\mathcal{O}_{Y_{\mathcal{J}}}\oplus\mathcal{J}\mathcal{O}_{Y}\)
(\ie the graded algebra associated with the filtration
\(\mathcal{O}_{Y}\supset\mathcal{J}\mathcal{O}_{Y}\)). Then \(L_{X}(Y)\) is
identified with \(X_{\mathcal{J}}(\widetilde{Y}_{\mathcal{J}})\) and
\(L'_{X}(Y)\) with the kernel of the morphism
\(p\colon X_{\mathcal{J}}(\widetilde{Y}_{\mathcal{J}})
\to X_{\mathcal{J}}(Y_{\mathcal{J}})\) induced by the ``zero section''
\(Y_{\mathcal{J}}\to\widetilde{Y}_{\mathcal{J}}\)
(\ie by the morphism of \(\mathcal{O}_{Y_{\mathcal{J}}}\)-algebras
\(\mathcal{O}_{\widetilde{Y}_{\mathcal{J}}}\to\mathcal{O}_{Y_{\mathcal{J}}}\)
vanishing on the ideal \(\mathcal{J}\mathcal{O}_{Y}\)). One therefore has,
for every \(Y\to S\), a split exact sequence of monoids, functorial in~\(Y\):
\[
\xymatrix{
1 \ar[r]
  & L'_{X}(Y) \ar[r]^{i}
  & L_{X}(Y) \ar@<0.5ex>[r]^{p}
  & X^{+}(Y) \ar@<0.5ex>[l]^{s} \ar[r]
  & 1\,.
}
\]

% typo: que loi -> que la loi
It remains to see that the monoid law induced on \(L'_{X}\) coincides with
its structure of abelian group. Let us denote by~\(\mu\) the monoid law
of~\(L_{X}\) and by~\(e\) its unit section; one must show that for all
\(m,m'\in L'_{X}(Y)\), one has
\[
\mu(m\cdot e,\,m'\cdot e)=(m+m')\cdot e.
\]
This can be seen in one or the other of the following ways. On the one hand,
one can take up again the proof of the equality~(0.1.10~\((\ast)\)) upon
replacing the morphism \(f\colon X\to W\) which figures there by the morphism
\(\mu\colon L_{X}\times_{S}L_{X}\to L_{X}\). Identifying
\(X^{+}(Y)=X_{\mathcal{J}}(Y_{\mathcal{J}})\) with its image by~\(s\) in
\(L_{X}(Y)=X_{\mathcal{J}}(\widetilde{Y}_{\mathcal{J}})\) one obtains that,
for all \(g,g'\in X_{\mathcal{J}}(Y_{\mathcal{J}})\) and
\(m,m'\in L'_{X}(Y)\), one has
\begin{equation*}
(\star)\qquad
\mu(m\cdot g,\,m'\cdot g')
=
L_{\mu}^{(g,g')}(m,m')\cdot\mu(g,g'),
\end{equation*}
where \(L_{\mu}^{(g,g')}\) denotes the derived morphism of~\(\mu\) at the
point~\((g,g')\) (\ie \(\widetilde{Y}_{\mathcal{J}}\) lies over
\(L_{X}\times_{S}L_{X}\) via~\((g,g')\)). In particular, one has
\(\mu(m\cdot e,\,m'\cdot e)=L_{\mu}^{(e,e)}(m,m')\cdot e\); but
\(L_{\mu}^{(e,e)}(m,m')=L_{\ell_{e}}(m')+L_{r_{e}}(m)\), where \(\ell_{e}\)
(\resp \(r_{e}\)) denotes the left translation (\resp the right translation)
by~\(e\), which is the identity map of~\(X_{\mathcal{J}}\), whence
\(L_{\mu}^{(e,e)}(m,m')=m+m'\).

Or else, one can proceed as follows (\cf\ the proof of \textbf{[DG70]},
\S~II.4, Th.~3.5). By lemma~0.1.11, the formation of \(X^{+}\) and of
\(L_{X}\) ``commutes with the product''
% typo: il en donc -> il en est donc
and the same therefore holds for~\(L'_{X}\); it follows that the morphism
\(\mu'\colon L'_{X}\times L'_{X}\to L'_{X}\), induced by~\(\mu\), is a
homomorphism for the structure of abelian groups. One then deduces from
lemma~3.10 of Exp.~II that \(\mu'\) coincides with the abelian group law.
\end{proof}

\numpara{0.8}%
{\renewcommand{\thefootnote}{(27)}\nde{One has added here the number~0.8 in
order to mark the return to the original.}}
Suppose now that \(X\) is an \(S\)-scheme such that \(X_{\mathcal{J}}\) is an
\(S_{\mathcal{J}}\)-group. Suppose that there exists an \(S\)-morphism
\[
P\colon X\times_{S}X \to X
\]
such that the morphism obtained by base change
\[
P_{\mathcal{J}}\colon
  X_{\mathcal{J}}\times_{S_{\mathcal{J}}}X_{\mathcal{J}}
  \to X_{\mathcal{J}}
\]
is the group law of~\(X_{\mathcal{J}}\). (An important particular case of the
preceding situation will be the case where \(X\) is an \(S\)-group and where
one takes for~\(P\) its group law.) One deduces from this a morphism
\[
L_{P}\colon
  L_{X}\times_{S}L_{X}
  \simeq
  L_{X\times_{S}X}
  \to
  L_{X}
\]
which, in fact, does not depend on~\(P\), for it is computed by means of the
group law \(P_{\mathcal{J}}\) of~\(X_{\mathcal{J}}\) as we are going to see
now.%
{\renewcommand{\thefootnote}{(28)}\nde{One has set out the original in detail
in what follows.}}
Indeed, by~(ii) and~(iii) of~0.2, for every \(Y\to S\) and
\(x,x'\in X(Y)\), \(m,m'\in L'_{X}(Y)\), one has
\[
P(m\cdot x,\,m'\cdot x')
=
P\bigl((m,m')\cdot(x,x')\bigr)
=
L_{P}^{(x,x')}(m,m')\cdot\mu(g,g')
\]
where \(g\) (\resp \(g'\)) is the image of~\(x\) (\resp \(x'\)) in
\(X^{+}(Y)\). Moreover (\cf\ the proof of~0.10), \(L_{P}^{(x,x')}\) equals
\(L_{\mu}^{(g,g')}\) and, by~0.7.bis~\((\star)\), the latter is the element
of \(L'_{X}(Y)\) defined by the following equality in \(L_{X}(Y)\):
\[
L_{\mu}^{(g,g')}(m,m')\cdot\mu(g,g')
=
\mu(m\cdot g,\,m'\cdot g'),
\]
that is, if one denotes by~\(\times\) (instead of~\(\mu\)) the group law
of~\(L_{X}\) and by~\(\Ad\) ``the adjoint operation'' of~\(X^{+}\) on
\(L'_{X}\) (which factors through~\(X_{0}\) and which is induced by the
adjoint operation of~\(X_{0}\) on \(\omega^{1}_{X_{0}/S_{0}}\)), one obtains
that
\[
L_{\mu}^{(g,g')}(m,m')\times g\times g'
=
m\times g\times m'\times g'
=
\bigl(m\times\Ad(g)(m')\bigr)\times g\times g'
\]
whence finally \(L_{P}^{(x,x')}(m,m')=m\times\Ad(g)(m')\). One therefore
obtains the following:

\begin{proposition}{0.8}
\label{III.0.8}
Let \(P\colon X\times_{S}X\to X\) be an \(S\)-morphism such that
\(P_{\mathcal{J}}\) equips \(X_{\mathcal{J}}\) with a structure of
\(S_{\mathcal{J}}\)-group. Let us denote by~\(\times\) the group law of
\(L'_{X}\) and by \((m,x)\mapsto m\cdot x\) the morphism
\(L'_{X}\times_{S}X\to X\) defining the action of \(L'_{X}\) on~\(X\), and let
\(\Ad\colon X^{+}\to\Aut_{S\text{-gr.}}(L'_{X})\) be ``the adjoint operation''
of~\(X^{+}\) on \(L'_{X}\) (which is induced by the adjoint operation of
\(X_{0}\) on \(\omega^{1}_{X_{0}/S_{0}}\)). Then, for every \(S'\to S\) and
\(x,x'\in X(S')\), \(m,m'\in L'_{X}(S')\), one has:
\begin{equation*}
(0.8.1)\qquad
P(m\cdot x,\,m'\cdot x')
=
\bigl(m\times\Ad\,p_{X}(x)(m')\bigr)\cdot P(x,x').
\end{equation*}
\end{proposition}

If \(X\) is an \(S\)-group, one will denote by~\(*\) its law, by~\(e\) its unit
section, and by~\(i\) the \(S\)-morphism defined by:
\[
i(m)=m\cdot e,
\]
for every \(S'\to S\) and \(m\in L'_{X}(S')\).

\begin{corollary}{0.9}
\label{III.0.9}
Let \(X\) be an \(S\)-group. Then \(X^{+}\) is equipped naturally with a
structure of \(S\)-group, and \(p_{X}\) is a morphism of \(S\)-groups.
Moreover, the \(S\)-morphism
\[
i\colon L'_{X}\to X,\qquad m\mapsto m\cdot e
\]
is an isomorphism of \(S\)-groups of \(L'_{X}\) onto \(\Ker(p_{X})\), and one
has, for all \(S'\to S\), \(x'\in X(S')\), \(m\in L'_{X}(S')\):
\begin{equation*}
(0.9.1)\qquad
m\cdot x'=(m\cdot e)*x'=i(m)*x'.
\end{equation*}
\end{corollary}

The first two assertions have already been proved in~0.1.2. As \(X\) is
formally principal homogeneous over~\(X^{+}\) under
\(L_{X}=L'_{X}\times_{S}X^{+}\), the morphism~\(i\) is indeed an isomorphism
of \(S\)-functors of \(L'_{X}\) onto the kernel of~\(p_{X}\). The fact that
\(i\) is a morphism of groups and the formula~(0.9.1) result from the
formula~(0.8.1) applied respectively to \(x=x'=e\), and to \(x=e\), \(m'=1\).

\begin{corollary}{0.10}
\label{III.0.10}
Let \(X\) be an \(S\)-group. With the preceding notations, for every \(S'\to S\)
and all \(x\in X(S')\) and \(m'\in L'_{X}(S')\), one has
\begin{equation*}
(0.10.1)\qquad
x*i(m')*x^{-1}
=
i\bigl(\Ad\,p_{X}(x)(m')\bigr).
\end{equation*}
\end{corollary}

This results from the equality \(i(m')*x^{-1}=m'\cdot x^{-1}\) and
from~(0.8.1) applied to \(m=1\) and \(x'=x^{-1}\).

When \(X\) is an \(S\)-group, we have therefore determined explicitly the
kernel of \(X\to X^{+}\) and the operation of the inner automorphisms of~\(X\)
on this kernel. We are now going to see that one can do likewise for certain
\(S\)-functors in groups, not necessarily representable. One case will be
useful to us, that of the functors \(\Aut\) (I~1.7).
Let us state at once:

\begin{proposition}{0.11}
\label{III.0.11}
Let \(E\) be an \(S\)-scheme. Set \(X=\Aut_{S}(E)\). The kernel of
the morphism of \(S\)-functors in groups
\[
p_{X}\colon X\to X^{+}
\]
is identified canonically with the commutative \(S\)-functor in groups
\(L'_{X}\) defined by
\[
\Hom_{S}(Y,L'_{X})
=
\Hom_{\mathcal{O}_{E_{0}\times_{S_{0}}Y_{0}}}
  \bigl(\Omega^{1}_{E_{0}/S_{0}}\otimes_{\mathcal{O}_{S_{0}}}\mathcal{O}_{Y_{0}},\,
  \mathcal{J}\mathcal{O}_{E\times_{S}Y}\bigr),
\]
where \(Y\) denotes a variable \(S\)-scheme.
\end{proposition}

Indeed, if \(Y\) is a variable \(S\)-scheme, one has
\(\Hom_{S}(Y,X)=\Aut_{Y}(E\times_{S}Y)\), and
\[
\Hom_{S}(Y,X^{+})
=
\Hom_{S}(Y_{\mathcal{J}},X)
=
\Aut_{Y_{\mathcal{J}}}(E\times_{S}Y_{\mathcal{J}})
=
\Aut_{Y_{\mathcal{J}}}\bigl((E\times_{S}Y)\times_{Y}Y_{\mathcal{J}}\bigr).
\]
By applying~0.7~b) to the \(Y\)-scheme \(E\times_{S}Y\), one obtains an
isomorphism of groups:
\[
\Hom_{S}(Y,L'_{X})
\simeq
\Ker\bigl(\Hom_{S}(Y,X)\to\Hom_{S}(Y,X^{+})\bigr),
\]
an isomorphism which one easily verifies to be functorial in the \(S\)-scheme
\(Y\). One therefore obtains an isomorphism of \(S\)-groups
\[
L'_{X}\simeq\Ker(X\to X^{+}),
\]
which completes the proof of proposition~0.11.

\begin{corollary}{0.12}
\label{III.0.12}
{\renewcommand{\thefootnote}{(29)}\nde{\normalfont One has changed ``Remark''
into ``Corollary''.}}
One keeps the notations of~0.11: \(E\) is an \(S\)-scheme and
\(X=\Aut_{S}(E)\). One has a natural operation~\(f\) of~\(X\) on \(L'_{X}\)
defined in the following manner. For every \(S\)-scheme~\(Y\), one has
\[
\Hom_{S}(Y,X)=\Aut_{Y}(E\times_{S}Y)
\]
and
\[
\Hom_{S}(Y,L'_{X})
=
\Hom_{\mathcal{O}_{E_{0}\times_{S_{0}}Y_{0}}}
  \bigl(\Omega^{1}_{E_{0}\times_{S_{0}}Y_{0}/Y_{0}},\,
  \mathcal{J}\mathcal{O}_{E\times_{S}Y}\bigr)
\]
(N.B.\ \(\Omega^{1}_{E_{0}/S_{0}}\otimes_{S_{0}}\mathcal{O}_{Y_{0}}
\simeq\Omega^{1}_{E_{0}\times_{S_{0}}Y_{0}/Y_{0}}\),
\cf\ EGA~IV\(_4\), 16.4.5); the first group operates on the second by
transport of structure and this operation is indeed functorial in~\(Y\).
One then has the formula:
\begin{equation*}
(0.12.1)\qquad
x\,i(m)\,x^{-1}=i\bigl(f(x)m\bigr),
\end{equation*}
for every \(Y\to S\) and all \(x\in\Hom_{S}(Y,X)\),
\(m\in\Hom_{S}(Y,L'_{X})\).
\end{corollary}

Indeed, this results from~0.7~c) applied to the \(Y\)-scheme
\(E\times_{S}Y\).

\begin{enoncerm}{Recall}{0.13}
\label{III.0.13}
The direct image of a quasi-coherent module by a morphism of finite
presentation is quasi-coherent. Under the same conditions, the formation of
the direct image commutes with \emph{flat} base change: in the situation
\[
\xymatrix@C=3.2em@R=1.6em{
T \ar[d]_{f}
  & T'=T\times_{S}S' \ar[l]_{g'} \ar[d]^{f'} \\
S
  & S' \ar[l]_{g}\,,
}
\]
if one supposes \(f\) (and hence \(f'\)) of finite presentation and \(g\) (and
hence \(g'\)) \emph{flat}, one has for every quasi-coherent
\(\mathcal{O}_{T}\)-module~\(\mathcal{F}\)
\[
f_{*}(\mathcal{F})\otimes_{\mathcal{O}_{S}}\mathcal{O}_{S'}
=
f'_{*}(\mathcal{F}\otimes_{\mathcal{O}_{S}}\mathcal{O}_{S'}),
\]
or, in a more aesthetic manner
\[
g^{*}(f_{*}(\mathcal{F}))=f'_{*}(g'^{*}(\mathcal{F})).
\]
These two facts are more generally valid for a morphism~\(f\) quasi-compact
and quasi-separated, \cf\ EGA~I, 9.2.1 and EGA~III\(_1\), 1.4.15 in the
quasi-compact separated case (taking into account EGA~III\(_2\),
Err\(_{\mathrm{III}}\)~25) and EGA~IV\(_1\), 1.7.4 and~1.7.21.
\end{enoncerm}

\begin{remark}{0.14}
\label{III.0.14}
{\renewcommand{\thefootnote}{(30)}\nde{One has set out the first part of this
remark in detail, and one has added a second part.}}
Let us take up again the notations of~0.11: let \(E\) be an \(S\)-scheme,
\(X=\Aut_{S}(E)\) and \(L'_{X}\) the commutative \(S\)-functor in groups defined
by:
\begin{align*}
L'_{X}(Y)
&=
\Hom_{\mathcal{O}_{E_{0}\times_{S_{0}}Y_{0}}}
  \bigl(\Omega^{1}_{E_{0}/S_{0}}
    \otimes_{\mathcal{O}_{S_{0}}}\mathcal{O}_{Y_{0}},\,
  \mathcal{J}\mathcal{O}_{E\times_{S}Y}\bigr) \\
&=
\Hom_{\mathcal{O}_{E\times_{S}Y}}
  \bigl(\Omega^{1}_{E/S}\otimes_{\mathcal{O}_{S}}\mathcal{O}_{Y},\,
  \mathcal{J}\mathcal{O}_{E\times_{S}Y}\bigr) \\
&=
\Gamma\bigl(E\times_{S}Y,\,
  \mathcal{H}\mathrm{om}_{\mathcal{O}_{E\times_{S}Y}}
    \bigl(\Omega^{1}_{E/S}\otimes_{\mathcal{O}_{S}}\mathcal{O}_{Y},\,
    \mathcal{J}\mathcal{O}_{E\times_{S}Y}\bigr)\bigr).
\end{align*}
\end{remark}
